\documentclass[10pt,a4paper]{article}
\usepackage[margin=1in]{geometry}
\usepackage{amsmath,amssymb}
\usepackage{graphicx}
\usepackage[hidelinks]{hyperref}
\usepackage{natbib}
\usepackage{booktabs}
\usepackage{enumitem}
\usepackage{orcidlink}

\title{WrinkleFE: An Open-Source Finite Element Package for Strength Prediction of Wrinkled Composite Laminates}

\author{Rani Elhajjar\textsuperscript{1}\,\orcidlink{0000-0003-4778-5186}}

\date{}

\begin{document}

\maketitle

\noindent\textsuperscript{1} Department of Civil and Environmental Engineering, College of Engineering and Applied Science, University of Wisconsin--Milwaukee, Milwaukee, WI 53201, USA

\bigskip

% ======================================================================
\section*{Summary}
% ======================================================================

WrinkleFE is a Python finite element package for predicting strength and stiffness knockdown in composite laminates containing fiber waviness defects.
The software combines Budiansky--Fleck kink-band theory \citep{BudianskyFleck1993}, Jin et al.\ Gaussian-sinusoidal wrinkle geometry parameterization \citep{Jin2026}, physics-based morphology factors derived from wavy plate theory \citep{Timoshenko1959}, the NASA LaRC04/05 failure criteria \citep{Pinho2005,DavilaCamanho2005}, and a three-mechanism tension knockdown model based on stress transformation and curved-beam interlaminar mechanics.
WrinkleFE supports five wrinkle morphologies---stack ($\varphi=0$), convex ($\varphi=\pi/2$), concave ($\varphi=-\pi/2$), uniform (constant through-thickness amplitude), and graded (amplitude decaying from midplane to surface with configurable decay floor)---providing comprehensive coverage of the defect configurations encountered in practice \citep{EllHajjar2023}.
The software provides both loading-specific analytical predictions (compression: Budiansky--Fleck; tension: fibre/matrix/out-of-plane interaction) and full three-dimensional hexahedral finite element analysis with the LaRC04/05 criterion.
A PyQt6 graphical interface enables interactive configuration of wrinkle geometry, material selection from a built-in database, mesh generation, and real-time visualization of stress fields, failure mechanisms, and strength and modulus retention factors.
The package is installable via \texttt{pip install wrinklefe}.

% ======================================================================
\section*{Statement of Need}
% ======================================================================

Fiber waviness is a common manufacturing defect in composite structures, arising from ply consolidation, tool--part interaction, automated fiber placement, resin flow during cure, and co-bonding processes \citep{Wisnom1990,EllHajjar2023}.
Even small fiber misalignment angles of 5--10$^\circ$ can reduce compression strength by 30--60\% through kink-band formation, while tension strength is also degraded through interlaminar stress concentration and matrix cracking at lower severity levels \citep{Hsiao1996,Elhajjar2025,Mukhopadhyay2015}.
When two wrinkles interact at adjacent ply interfaces---the dual-wrinkle configuration studied by \citet{Jin2026}---the relative phase offset between wrinkles creates distinct morphologies (stack, convex, concave) with markedly different failure behaviors.
Through-thickness amplitude variation further distinguishes uniform waviness (constant amplitude across all plies) from graded waviness (amplitude decaying from the wrinkle center toward the surface), as commonly observed in manufactured parts \citep{EllHajjar2023}.

No open-source software currently integrates wrinkle geometry parameterization, morphology-dependent interlaminar damage, physics-based failure criteria, and finite element stress analysis in a single framework.
WrinkleFE addresses this gap by providing a validated, modular, and extensible Python package that researchers and engineers can use directly or build upon.

% ======================================================================
\section*{State of the Field}
% ======================================================================

Closed-form analytical models based on Budiansky--Fleck kink-band theory \citep{BudianskyFleck1993} predict compression knockdown from a single misalignment angle but do not capture morphology effects, interlaminar damage, or tension failure mechanisms.
Existing finite element studies of wrinkled laminates \citep{Jin2026,Mukhopadhyay2015} provide high-fidelity damage predictions but rely on proprietary software with custom subroutine development, limiting reproducibility and transferability.

An open question in the field concerns the relative importance of compression versus tension knockdown.
While compression failure via kink-band instability has received extensive study \citep{BudianskyFleck1993,Fleck1997,Pimenta2015,Hsiao1996}, tension failure in wrinkled laminates has received less attention.
The assumption that tension knockdown is negligible is contradicted by experimental data showing 23--53\% strength reduction at moderate-to-severe waviness \citep{Elhajjar2025} and progressive delamination-to-fibre-fracture failure sequences observed in quasi-isotropic laminates \citep{Mukhopadhyay2015}.
WrinkleFE addresses this by implementing a physics-based three-mechanism tension model alongside the established compression framework.

WrinkleFE combines the speed of analytical models with the fidelity of finite element analysis in an open-source package.
The software implements the LaRC04/05 failure criterion \citep{Pinho2005}, which accounts for fibre kinking through misalignment-frame stress rotation---the key mechanism connecting wrinkle geometry to compression failure.
In-situ ply strengths use the thick-ply correction from \citet{Camanho2006}: $S_{12}^{is} = \sqrt{8G_{IIc}/(\pi\,t_{\text{eff}}\,\Lambda_{22}^0)}$, where $t_{\text{eff}}$ is the adjacent 0$^\circ$-ply group thickness, providing physically accurate matrix-mode predictions without requiring separate calibration.
WrinkleFE automates the entire workflow from wrinkle parameter specification through mesh generation, stress analysis, and failure assessment, with built-in comparison between wrinkled and pristine (baseline) laminates to produce physically meaningful retention factors.

% ======================================================================
\section*{Software Design}
% ======================================================================

WrinkleFE is organized into five core modules with well-defined interfaces:

\begin{description}[leftmargin=1.5em,labelindent=0em]
\item[Core] provides material property databases (four built-in carbon/epoxy systems: AS4/3501-6, IM7/8552 [default], T300/914, T700/2510), classical lamination theory, structured hexahedral mesh generation with wrinkle deformation, and the wrinkle geometry model.
The Gaussian-sinusoidal profile $z(x) = A \exp(-(x-x_0)^2/w^2)\cos(2\pi(x-x_0)/\lambda)$ from \citet{Jin2026} is centered in the specimen domain.
Five morphologies describe how wrinkle amplitude varies through the laminate thickness:
\begin{itemize}[nosep]
\item \textit{Stack} ($\varphi=0$): upper and lower wrinkle peaks aligned---produces symmetric interlaminar gap.
\item \textit{Convex} ($\varphi=\pi/2$): interface bulges outward---best case for compression (stabilizes against kink-band).
\item \textit{Concave} ($\varphi=-\pi/2$): interface pinches inward---worst case for compression (amplifies kink-band).
\item \textit{Uniform}: constant wrinkle amplitude across all plies (no through-thickness decay).
\item \textit{Graded}: amplitude decays linearly from midplane to surface with a configurable decay floor parameter ($0$--$1$), where $0$ gives full decay (zero amplitude at surface) and $1$ gives uniform amplitude.
\end{itemize}
Dual-wrinkle morphologies (stack, convex, concave) are realized by applying a geometric phase offset $\Delta x = \varphi\lambda/(2\pi)$ to the second wrinkle's profile evaluation during mesh deformation.
The morphology factor $M_f(\varphi) = \exp(-\alpha_{\text{asym}}\sin\varphi - \alpha_{\text{offset}}(1-|\cos\varphi|))$ modifies the effective fiber angle according to loading mode (compression: $\alpha_{\text{asym}}=0.288$; tension: $\alpha_{\text{asym}}=0.033$, $\alpha_{\text{offset}}=0.183$).

\item[Solver] implements a displacement-based static solver with 8-node hexahedral elements (full $2\times2\times2$ Gauss integration), direct sparse factorization, and compression or tension boundary conditions via prescribed displacement.

\item[Failure] implements two loading-specific analytical models and the LaRC04/05 criterion \citep{Pinho2005,DavilaCamanho2005} evaluated at every Gauss point.

\textit{Compression analytical model:} CLT-weighted Budiansky--Fleck kink-band theory, $\text{KD}_{\text{lam}} = f_0 \cdot 1/(1+\theta_{\text{eff}}/\gamma_Y^{eff}) + (1-f_0)$, where $f_0$ is the CLT axial stiffness fraction of $0^\circ$ plies and $\gamma_Y^{eff}$ is the layup-dependent effective yield strain computed from ply confinement: $\gamma_Y^{eff} = 0.020 + 0.050\,f_{\text{confined}}$.
Here $f_{\text{confined}}$ is a weighted confinement fraction: each $0^\circ$ ply scores 1.0 if both neighbors are off-axis, 0.5 if one neighbor is off-axis, and 0.0 if both neighbors are $0^\circ$ (block interior).
The CLT weighting separates two physical effects: confinement (captured by $\gamma_Y^{eff}$) and load redistribution to off-axis plies (captured by $f_0$), enabling the model to handle both dispersed and blocked layups \citep{Fleck1997,Pimenta2015}.

\textit{Tension analytical model:} A physics-based three-mechanism model determines the 0$^\circ$-ply knockdown as the minimum of three competing failure modes, weighted by CLT axial stiffness fractions:
(1)~fibre tension $\cos^2\theta$ from load projection \citep{Pinho2005};
(2)~matrix tension via Hashin interaction of $\sigma_{22} = \sigma\sin^2\theta$ and $\tau_{12} = \sigma\sin\theta\cos\theta$ with in-situ strengths: $Y_T^{is} = 1.12\sqrt{2}\,Y_T$ (Pinho Eq.~47) and $S_{12}^{is}$ from the thick-ply formula $S_{12}^{is} = \sqrt{8G_{IIc}/(\pi\,t_{\text{eff}}\,\Lambda_{22}^0)}$ \citep{Camanho2006}, where $t_{\text{eff}} = n_{\text{adj}} t_{\text{ply}}$ is the adjacent 0$^\circ$-ply group thickness and $G_{IIc} \approx 1.0$~N/mm for typical carbon/epoxy;
(3)~out-of-plane delamination from curved-beam interlaminar stresses at the wrinkle crest, where $\sigma_{33} = X_T \cdot h_{\text{eff}} \cdot \kappa_{\max}$ with $h_{\text{eff}} = n_{\text{adj}} t_{\text{ply}}$ and $\kappa_{\max} = (2\pi/\lambda)^2 A$ (peak curvature).
The delamination knockdown uses a smooth progressive-damage interaction: $\text{KD}_{\text{oop}} = 1/\sqrt{1 + (\sigma_{33}/Y_T)^2}$.
The laminate knockdown is $\text{KD}_{\text{lam}} = f_0 \cdot \min(\text{KD}_{\text{fiber}}, \text{KD}_{\text{matrix}}, \text{KD}_{\text{oop}}) + (1-f_0)$.
For graded morphology, the knockdown is averaged over all 0$^\circ$ plies at their local through-thickness angles $\theta_p = \theta_{\max} \cdot B_p$.

\textit{LaRC04/05 FE criterion:}
Fibre kinking ($\sigma_1 < 0$) rotates stresses into the misalignment frame by $\varphi = \varphi_0 + \varphi_c$.
Fibre tension ($\sigma_1 \geq 0$): $\text{FI} = (\sigma_1/X_T)^2 + (\tau_{12}/S_{12})^2$.
Matrix failure uses Mohr--Coulomb fracture-plane search with in-situ strengths.
Boundary elements are excluded from failure assessment to remove artificial stress concentrations at constrained faces.

\item[Analysis] orchestrates the pipeline: configuration $\to$ laminate $\to$ wrinkle $\to$ mesh $\to$ solve $\to$ failure $\to$ retention factor.
The retention factor is computed by running a second baseline analysis with zero wrinkle amplitude and reporting the ratio of pristine to wrinkled failure indices, providing a normalized strength metric independent of applied load magnitude.
A modulus retention factor $E_{\text{wrinkled}}/E_{\text{pristine}}$ is also extracted from the FE solution by comparing mean axial stresses under the same applied strain.

\item[GUI] offers a PyQt6 interface with interactive controls for material selection, layup specification, wrinkle geometry, and mesh density.
Three visualization tabs display wrinkle profiles, mesh cross-sections, and stress contours with a mechanism decomposition bar chart showing strength and modulus retention factors.
The bar chart annotation adapts to the loading mode, displaying the controlling failure mechanism and the confinement-derived $\gamma_Y^{eff}$.
Analysis runs in a background thread with an estimated-time countdown.
\end{description}

The package includes over 400 unit and integration tests covering all modules.
A command-line interface supports batch analysis, morphology comparison, and parametric sweeps.

% ======================================================================
\section*{Research Impact}
% ======================================================================

WrinkleFE has been validated against 26 experimental data points from two independent datasets: \citet{Elhajjar2025} (T700GC/2510, dispersed layup, uniform waviness, 20 points) and \citet{Mukhopadhyay2015} (IM7/8552, blocked layup, graded waviness, 6 points), achieving 24/26 PASS within $\pm 20\%$ tolerance (overall MAE = 9.9\%).
The CLT-weighted confinement model $\gamma_Y^{eff} = 0.020 + 0.050\,f_{\text{confined}}$ automatically computes the effective yield strain from the layup stacking sequence, separating matrix confinement from load redistribution and eliminating the need for per-configuration calibration.

\textit{Compression:} The CLT-weighted Budiansky--Fleck model matches 12/13 data points from \citet{Elhajjar2025} (D/T = 0.003--0.30, MAE = 9.5\%) within the $\pm 20\%$ tolerance.
The confinement model gives $\gamma_Y^{eff} = 0.062$ for the Elhajjar layup ($f_{\text{confined}} = 0.833$, $f_0 = 0.714$).
The single case that exceeds the tolerance (D/T = 0.300) occurs where the CLT floor $(1-f_0)$ saturates---at this extreme severity level, the off-axis plies also begin to degrade, which the current model does not capture.
The wrinkle wavelength is estimated as $\lambda = 19.9 A$ (with $\lambda_{min} = 8.2$~mm), calibrated from the micrograph geometry reported in \citet{Elhajjar2014}.

\textit{Tension:} The three-mechanism analytical model matches 7/7 data points from \citet{Elhajjar2025} within $\pm 20\%$ (MAE = 6.2\%).
The thick-ply in-situ shear strength correction \citep{Camanho2006} is essential for accurate matrix-mode predictions; the thin-ply approximation ($S_{12}^{is} = \sqrt{2}\,S_{12}$) underestimates the in-situ strength for adjacent 0$^\circ$-ply groups, causing premature matrix-dominated knockdown.
The model captures three distinct failure regimes: (1)~at low D/T ($< 0.05$), curved-beam $\sigma_{33}$ from wrinkle curvature drives delamination; (2)~at moderate D/T ($0.05$--$0.15$), out-of-plane and matrix mechanisms compete; (3)~at high D/T ($> 0.15$), Hashin matrix interaction dominates.
These tension knockdown levels---23--53\% for moderate-to-severe waviness---are consistent with experimental observations by both \citet{Elhajjar2025} and \citet{Mukhopadhyay2015}, contradicting the common assumption that tension knockdown is negligible.

\textit{Graded waviness:} For graded wrinkles where amplitude decays from the midplane to the surface, the analytical model averages the knockdown over all 0$^\circ$ plies at their local angles $\theta_p = \theta_{\max} \cdot B_p$, where $B_p$ is the through-thickness grading factor.
This captures the physical reality that surface plies see minimal waviness while midplane plies see the full amplitude.
With graded averaging, the model matches 2/3 compression and 3/3 tension (vs.\ ultimate failure) data points from \citet{Mukhopadhyay2015} (IM7/8552, $[{+}45_2/90_2/{-}45_2/0_2]_{3s}$, MAE = 15.0\%).
Notably, \citet{Mukhopadhyay2015} report both a first load-drop (delamination initiation) and final failure (fibre fracture) for the higher-severity tension specimens.
Since the controlling mechanisms in the tension model---curved-beam $\sigma_{33}$ and Hashin matrix interaction---physically predict damage onset rather than ultimate collapse, the model predictions fall between these two experimental values (Figure~\ref{fig:mukhopadhyay}), consistent with the progressive damage sequence observed experimentally.

Overall, the analytical model achieves 24/26 PASS across both datasets (MAE = 9.9\%).
Both the compression and tension models are physics-based: material properties (elastic constants and ply-level strengths) and the layup stacking sequence are the only inputs required.

\begin{figure}[htbp]
\centering
\includegraphics[width=\textwidth]{fig_validation_elhajjar.png}
\caption{Validation against \citet{Elhajjar2025} (T700/2510, uniform waviness).  Red: compression (CLT-weighted Budiansky--Fleck).  Blue: tension (three-mechanism model with thick-ply in-situ correction).  19/20 within $\pm 20\%$, MAE = 8.4\%.}
\label{fig:validation}
\end{figure}

\begin{figure}[htbp]
\centering
\includegraphics[width=0.85\textwidth]{fig_validation_mukhopadhyay.png}
\caption{Validation against \citet{Mukhopadhyay2015} (IM7/8552, graded waviness).  Green circles: compression (catastrophic failure).  Purple squares: tension ultimate (fibre fracture).  Purple diamonds: tension first load-drop (delamination onset).  Model curves use graded ply-averaging.  The tension model predicts between initiation and ultimate failure, consistent with its onset-based mechanisms.}
\label{fig:mukhopadhyay}
\end{figure}

% ======================================================================
\section*{Scope and Limitations}
% ======================================================================

The following limitations should be considered when interpreting WrinkleFE predictions and when applying the software to configurations beyond those validated here.

\textit{What is validated:}
The model is validated against 26 data points from two datasets: \citet{Elhajjar2025} (T700GC/2510, dispersed layup, uniform waviness, 19/20 PASS) and \citet{Mukhopadhyay2015} (IM7/8552, blocked layup, graded waviness, 5/6 PASS).
The CLT-weighted confinement model ($\gamma_Y^{eff} = 0.020 + 0.050\,f_{\text{confined}}$) is calibrated from the Elhajjar dataset plus the UD baseline.
The confinement coefficient ($\alpha = 0.050$) may need adjustment for significantly different fibre/matrix combinations (e.g., thermoplastic matrices, glass fibres).
For graded wrinkles, the analytical model averages the knockdown over all 0$^\circ$ plies at their local through-thickness angles, which substantially improves accuracy compared to using the peak angle alone.

\textit{What is assumed:}
The model assumes linear elastic material behavior, quasi-static loading, and small deformations.
Progressive damage evolution (matrix cracking, delamination growth, fibre breakage) is not modeled---the failure criteria identify onset, not propagation.
The wrinkle geometry is idealized as a Gaussian-enveloped sinusoid; real manufacturing defects may exhibit asymmetric, non-periodic, or three-dimensional out-of-plane waviness patterns \citep{EllHajjar2023}.
Interaction between waviness and other defect types (voids, porosity, ply drops, foreign object inclusions) is not captured.
Environmental effects (temperature, moisture) and fatigue loading are outside the current scope.

\textit{Generalizability:}
The compression model (Budiansky--Fleck) rests on well-established mechanics that applies broadly to fibre-dominated composites \citep{Fleck1997}.
The tension model's three-mechanism structure---fibre projection, matrix interaction, and curved-beam interlaminar stress---uses general mechanics principles but has been validated only for carbon/epoxy quasi-isotropic layups.
The tension model uses the thick-ply in-situ shear strength correction $S_{12}^{is} = \sqrt{8G_{IIc}/(\pi\,t_{\text{eff}}\,\Lambda_{22}^0)}$ from \citet{Camanho2006} with a typical $G_{IIc} \approx 1.0$~N/mm; users with material-specific fracture toughness data may substitute their own value for improved accuracy.
For layups with few or no 0$^\circ$ plies, the CLT-weighted formulation would require revisiting, and the adjacent-ply-group thickness $h_{\text{eff}}$ would change.
The morphology factor (stack, convex, concave) derives from the phase offset between two wrinkles at adjacent interfaces and would need extension for configurations with more than two interacting wrinkles or wrinkles spanning multiple ply interfaces.

% ======================================================================
\section*{AI Usage Disclosure}
% ======================================================================

Anthropic's Claude Code was used as a development assistant during the creation of WrinkleFE.
The AI tool assisted with Python code generation, unit test writing, documentation drafting, and GUI implementation.
All scientific methodology---including the selection of failure criteria, the derivation of the morphology factor from wavy plate theory, the wrinkle geometry parameterization, and the validation strategy against published data---was directed by the author based on established composite mechanics theory and the referenced publications.
The author reviewed, tested, and validated all AI-generated code before inclusion.

% ======================================================================
\section*{Acknowledgments}
% ======================================================================

The author thanks the developers of NumPy, SciPy, matplotlib, and PyQt6 for the open-source scientific computing ecosystem upon which WrinkleFE is built.

% ======================================================================
\bibliographystyle{plainnat}
\bibliography{paper}

\end{document}
